\section{案例推演：企业级 Agent 平台的分阶段建设路线}

\subsection{阶段 0：先做 ``可控的最小系统''}
很多团队上来就尝试多 Agent、复杂工具编排和自动化闭环，结果在第一周就被日志噪声、链路超时和错误重试拖垮。更务实的路径是先构建一个可控最小系统（Minimum Controllable System）：
\begin{itemize}
    \item 单 Agent + 单业务域 + 可回放日志；
    \item 检索源有限且可标注；
    \item 失败默认安全（fail-safe）而不是强行继续执行。
\end{itemize}

\begin{importantbox}{最小系统的目标不是 ``能力上限''，而是 ``行为确定性''}
只有当系统在固定输入下能稳定复现，后续优化才有意义。否则每次改动都无法判断收益来自模型提升还是偶然波动。
\end{importantbox}

\subsection{阶段 1：把检索与执行解耦}
课程里强调 RAG 与 Agent 并非替代关系。在工程实现上，建议把 ``知识检索'' 与 ``动作执行'' 解耦为两个独立层：
\begin{enumerate}
    \item 检索层负责召回候选事实与证据；
    \item 决策层负责选择工具与下一步动作；
    \item 执行层负责副作用操作并返回结构化结果。
\end{enumerate}

解耦后的直接收益是：可以分别评估召回质量、规划质量和执行成功率，避免把所有问题都归因给 ``模型幻觉''。这也更符合讲者提出的系统化思路。

\begin{knowledgebox}{推荐的最小观测指标}
\begin{itemize}
    \item 检索召回率（Top-k Evidence Recall）；
    \item 规划一致性（同问题多次执行的动作序列偏差）；
    \item 工具成功率（成功调用比例与平均重试次数）；
    \item 端到端任务成功率（按业务验收标准定义）。
\end{itemize}
\end{knowledgebox}

\subsection{阶段 2：引入流控与恢复策略}
在 48 分钟案例里，系统恢复失败本质是 ``恢复流量 > 服务恢复速度''。因此平台第二阶段的关键不是再堆功能，而是把流控和恢复设计前置。一个可执行的方案是：
\begin{itemize}
    \item 对每类任务设定并发预算（concurrency budget）；
    \item 对模型服务设定暖启动窗口（warm-up window）；
    \item 对下游工具调用设定熔断阈值（circuit breaker）；
    \item 对任务状态设定幂等恢复点（idempotent checkpoint）。
\end{itemize}

\begin{warningbox}{没有恢复策略的自动化等于放大器}
自动化链路会把单次错误放大到批量错误。对于企业场景， ``先恢复、再优化'' 比 ``先加功能、再救火'' 的总成本更低。
\end{warningbox}

\subsection{阶段 3：治理与合规上线}
当系统进入企业生产后，关注点会从 ``能不能做'' 转向 ``能不能长期可控地做''。这要求引入治理层：
\begin{itemize}
    \item 权限与审计：谁触发了什么动作，证据链是否完整；
    \item 数据边界：私有数据是否跨域泄漏；
    \item 合规约束：是否满足行业监管与留痕要求；
    \item 运营指标：成本、成功率、时延是否可持续。
\end{itemize}

\begin{table}[H]
\centering
\begin{tabular}{p{2.4cm}p{3.5cm}p{3.6cm}p{3.6cm}}
\toprule
\textbf{建设阶段} & \textbf{核心目标} & \textbf{必须交付物} & \textbf{退出条件} \\
\midrule
阶段 0 & 行为可复现 & 结构化日志、回放脚本、基线任务集 & 同输入结果稳定且可解释 \\
阶段 1 & 模块可诊断 & 检索/规划/执行分层指标面板 & 能定位性能与质量瓶颈来源 \\
阶段 2 & 故障可恢复 & 流控、熔断、幂等恢复点 & 峰值下无级联崩溃 \\
阶段 3 & 生产可治理 & 权限审计、数据边界、合规报告 & 满足业务与监管上线标准 \\
\bottomrule
\end{tabular}
\caption{企业级 Agent 平台从 PoC 到生产的四阶段路线图}
\end{table}

\subsection{本章小结}
平台建设应该遵循 ``可控性先于复杂度'' 的路径。先把复现、诊断、恢复、治理补齐，再追求更高自动化水平，才能把 Agent 变成长期可运营能力。

